\part*{Abstract}

Abundance and density data of unmarked animal populations are often estimated with camera trap
distance sampling.
Fundamentally, this relies on obtaining animal-to-camera distance data for each observation,
involving slow and inefficient methods of manual and subjective distance estimation.
This study investigates a novel and virtually unexplored approach to abundance/density estimation,
where a monocular depth estimation (MDE) pipeline is applied to a camera trap dataset to obtain detection
distance estimates which are subsequently used to model the density and abundance of the local population.
Estimates derived using two MDE models (DPT and Depth Anything) with two pixel depth sampling methods are
obtained, analysing all distance, density and abundance estimates from each workflow configuration.
The results show that calibrated DPT outperforms Depth Anything for both distance and
density/abundance estimation when applied to this dataset; however, both models generally overestimate
distance and underestimate density/abundance with respect to the manual approach.
To achieve more accurate density/abundance estimates with this pipeline, this study finds that the number
of missed detections across all distances must be reduced.
It is concluded that this study demonstrates proof-of-concept that MDE integrated camera trap distance
sampling can be used, in this context, to achieve estimates of density and abundance that approach
the accuracy of conventional means.
